\documentclass[compress]{beamer}
%\usepackage{graphicx}
\usepackage{epsfig}
\usepackage{pstricks,pst-plot}
\usepackage{mathpazo}
\usepackage{amsmath}
\usepackage{amsfonts}
\usepackage{amsmath,amssymb,amsthm,enumerate,graphicx}
\usepackage{ifthen,latexsym,syntonly}
\usepackage{cases}
\usepackage{setspace}
\usepackage{mathabx}
\usepackage{color}
\usepackage{natbib}
\usepackage{float}
\usepackage{epstopdf}
\usepackage{mathabx}
\usepackage{color}
\usepackage{calc,xspace}
\usepackage{bibentry}
\usepackage{multirow}
\usepackage{graphicx}
%\usepackage{subcaption}
%\usepackage[stable]{footmisc}
\setcounter{MaxMatrixCols}{10}
\usepackage{soul}
\usepackage{comment}
\usepackage{threeparttable}
\usepackage{booktabs}
\usepackage{makecell}
\usepackage{subfig}
\makeatletter
\let\@@magyar@captionfix\relax
\makeatother
\newcommand{\E}{\mathbb{E}}
\newcommand{\ept}{{\rm E}}
\newcommand{\ep}{{\rm e}}
\newcommand{\be}{\begin{equation}}
\newcommand{\ee}{\end{equation}}
\newcommand{\bq}{\begin{eqnarray}}
\newcommand{\eq}{\end{eqnarray}}
\newcommand{\bi}{\begin{itemize}}
\newcommand{\ei}{\end{itemize}}
\newcommand{\bc}{\begin{center}}
\newcommand{\ec}{\end{center}}
\renewcommand{\i}{\item}
\newcommand{\bnm}{\begin{enumerate}}
\newcommand{\enm}{\end{enumerate}}
\newcommand{\bfr}{\begin{frame}[allowframebreaks]}
\newcommand{\efr}{\end{frame}}
\newcommand{\fmt}[1]{\frametitle{\bc{#1}\ec}}
\newcommand{\Tr}{{\mathrm Tr}}
\newtheorem{remark}{Remark}

\usetheme{Hannover}
%\usetheme{Berkeley}
%\usetheme{Bergen}
%\usetheme{Boadilla}
%\usetheme{Singapore}
%\usetheme{Malmoe}
\usecolortheme[rgb={0,0,1}]{structure}
\usepackage[english]{babel}
\usepackage[latin1]{inputenc}
\usepackage{natbib}
\bibliographystyle{plainnat}
\newcommand{\JS}[1]{{\color{red}{#1}}}
\newcommand{\PH}[1]{{\color{red}{#1}}}

%\bibliographystyle{plainnat}
\setbeamercovered{transparent}
\setbeamertemplate{navigation symbols}{}
\setbeamertemplate {footline}
{\quad\strut\insertframenumber/\inserttotalframenumber\strut\quad}

%\setbeamertemplate {footline}{\quad\strut\insertframenumber}{\inserttotalframenumber\strut\quad}
\title[Climate Finance]{Climate Finance}
\author[Jos\'e Scheinkman]{Jos\'e A. Scheinkman}
\institute[Columbia University]{Columbia University}

\date[Climate Finance]{PhD course, Spring 2026}
\setbeamertemplate{navigation symbols}{}



  
%\setbeamertemplate{navigation symbols}{}
\setbeamertemplate {footline}
{\quad\strut\insertframenumber/\inserttotalframenumber\strut\quad}

\begin{document}
\maketitle
\section{Results}




\begin{frame}\fmt{Parameter uncertainty}
\bi \i Hamiltonian Monte Carlo 
\i 1043 sites
\i $\xi=1$ 
\i $P^{ee}=4.7$
\i $P^{ee}$ monotone increasing (decreasing) on $\xi$ (ambiguity aversion)
\ei

NEXT TWO SLIDES ARE FROM FIGURE 8
\end{frame}

\begin{frame}{Ambiguity adjustment, $b=\$0,$ $P^a =P^{s}$}
\begin{figure}
\centering
\begin{minipage}{.5\textwidth} \centering
\includegraphics[width=\linewidth]{figures/theta_distribution_916_b0.png}
\end{minipage}%
\begin{minipage}{.5\textwidth}  \centering
\includegraphics[width=\linewidth]{figures/gamma_distribution_817_b0.png}
\end{minipage}
\end{figure}
\end{frame}


\begin{frame}{Ambiguity adjustment, $b=\$15,$ $P^a =P^{s}$}
\begin{figure}
\centering
\begin{minipage}{.5\textwidth} \centering
\includegraphics[width=\linewidth]{figures/theta_distribution_1021_b15.png}
\end{minipage}%
\begin{minipage}{.5\textwidth}  \centering
\includegraphics[width=\linewidth]{figures/gamma_distribution_938_b15.png}
\end{minipage}
\end{figure}
\end{frame}

\begin{comment}
\begin{frame}{Ambiguity adjustment, $b=15,$ $P^a =P^{s}$}
\begin{figure}
\centering
\begin{minipage}{.5\textwidth} \centering
\includegraphics[width=\linewidth]{figures/aux_input/theta_distribution_78.png}
\end{minipage}%
\begin{minipage}{.5\textwidth}  \centering
\includegraphics[width=\linewidth]{figures/aux_input/gamma_distribution_78.png}
\end{minipage}
\end{figure}
\end{frame}
\end{comment}

\begin{frame}{Evolution of agricultural area: productivity ambiguity, $P^a =P^{s}$}

\begin{figure}[H]
\centering
\begin{minipage}{.5\textwidth} 
    \centering
    \includegraphics[width=1.1\linewidth]{figures/aggregate_percentage_Z_b0_pehmc_4.5_pedet_6.6.png}
    \vspace{-10pt}
    \caption*{(a) b=\$0}
\end{minipage}%
\begin{minipage}{.5\textwidth}
    \centering
    \includegraphics[width=1.1\linewidth]{figures/aggregate_percentage_Z_b15_pehmc_19.5_pedet_21.6.png}
    \vspace{-10pt}
    \caption*{(b) b=\$15}
\end{minipage}
\vspace{-10pt}
%\caption{Evolution of agricultural area under ambiguity neutrality and ambiguity aversion, using the corresponding shadow prices.  }
\label{aggz-hmc-1043}
\end{figure}

\end{frame}




\begin{comment}
\begin{frame}{Land allocation dynamics: productivity ambiguity}
\transduration{1.0}
\begin{figure}[H]
\centering
\begin{minipage}{\linewidth} % choose width suitably
\includegraphics[width=9.5cm]{figures/aux_input/z_2017_gridless.png}
\end{minipage}

\end{figure}
\end{frame}
\begin{frame}{Land allocation dynamics: productivity ambiguity}
\transduration{1}

\begin{figure}
\centering
\begin{minipage}{.5\textwidth} \centering
  \includegraphics[width=\linewidth]{figures/aux_input/map_z2018_25Sites_xinoAmbiguity_pe20_delta0_02.png}
\end{minipage}%
\begin{minipage}{.5\textwidth}  \centering
          \includegraphics[width=\linewidth]{figures/aux_input/map_z2018_25Sites_xi0.35_pe20_delta0_02.png}
\end{minipage}
\end{figure}
\end{frame}
\begin{frame}{Land allocation dynamics: productivity ambiguity}
\transduration{1}
\begin{figure}
\centering
\begin{minipage}{.5\textwidth} \centering
  \includegraphics[width=\linewidth]{figures/aux_input/map_z2019_25Sites_xinoAmbiguity_pe20_delta0_02.png}
\end{minipage}%
\begin{minipage}{.5\textwidth}  \centering
          \includegraphics[width=\linewidth]{figures/aux_input/map_z2019_25Sites_xi0.35_pe20_delta0_02.png}
\end{minipage}
\end{figure} 
\end{frame}

\begin{frame}{Land allocation dynamics: productivity ambiguity}
\transduration{1}

    \begin{figure}
\centering
\begin{minipage}{.5\textwidth} \centering
  \includegraphics[width=\linewidth]{figures/aux_input/map_z2020_25Sites_xinoAmbiguity_pe20_delta0_02.png}
\end{minipage}%
\begin{minipage}{.5\textwidth}  \centering
          \includegraphics[width=\linewidth]{figures/aux_input/map_z2020_25Sites_xi0.35_pe20_delta0_02.png}
\end{minipage}
\end{figure} 
\end{frame}

\begin{frame}{Land allocation dynamics: productivity ambiguity}
\transduration{1}

   \begin{figure}
\centering
\begin{minipage}{.5\textwidth} \centering
  \includegraphics[width=\linewidth]{figures/aux_input/map_z2021_25Sites_xinoAmbiguity_pe20_delta0_02.png}
\end{minipage}%
\begin{minipage}{.5\textwidth}  \centering
          \includegraphics[width=\linewidth]{figures/aux_input/map_z2021_25Sites_xi0.35_pe20_delta0_02.png}
\end{minipage}
\end{figure} 
\end{frame}

\begin{frame}{Land allocation dynamics: productivity ambiguity}
\transduration{1}
\begin{figure}
\centering
\begin{minipage}{.5\textwidth} \centering
  \includegraphics[width=\linewidth]{figures/aux_input/map_z2022_25Sites_xinoAmbiguity_pe20_delta0_02.png}
\end{minipage}%
\begin{minipage}{.5\textwidth}  \centering
          \includegraphics[width=\linewidth]{figures/aux_input/map_z2022_25Sites_xi0.35_pe20_delta0_02.png}
\end{minipage}
\end{figure}
\end{frame}

\begin{frame}{Land allocation dynamics: productivity ambiguity}
\transduration{1}
\begin{figure}
\centering
\begin{minipage}{.5\textwidth} \centering
  \includegraphics[width=\linewidth]{figures/aux_input/map_z2023_25Sites_xinoAmbiguity_pe20_delta0_02.png}
\end{minipage}%
\begin{minipage}{.5\textwidth}  \centering
          \includegraphics[width=\linewidth]{figures/aux_input/map_z2023_25Sites_xi0.35_pe20_delta0_02.png}
\end{minipage}
\end{figure}
\end{frame}

\begin{frame}{Land allocation dynamics: productivity ambiguity}
\transduration{1}

   \begin{figure}
\centering
\begin{minipage}{.5\textwidth} \centering
  \includegraphics[width=\linewidth]{figures/aux_input/map_z2024_25Sites_xinoAmbiguity_pe20_delta0_02.png}
\end{minipage}%
\begin{minipage}{.5\textwidth}  \centering
          \includegraphics[width=\linewidth]{figures/aux_input/map_z2024_25Sites_xi0.35_pe20_delta0_02.png}
\end{minipage}
\end{figure}
\end{frame}

\begin{frame}{Land allocation dynamics: productivity ambiguity}
\transduration{1}

    \begin{figure}
\centering
\begin{minipage}{.5\textwidth} \centering
  \includegraphics[width=\linewidth]{figures/aux_input/map_z2025_25Sites_xinoAmbiguity_pe20_delta0_02.png}
\end{minipage}%
\begin{minipage}{.5\textwidth}  \centering
          \includegraphics[width=\linewidth]{figures/aux_input/map_z2025_25Sites_xi0.35_pe20_delta0_02.png}
\end{minipage}
\end{figure}
\end{frame}

\begin{frame}{Land allocation dynamics: productivity ambiguity}
\transduration{1}

    \begin{figure}
\centering
\begin{minipage}{.5\textwidth} \centering
  \includegraphics[width=\linewidth]{figures/aux_input/map_z2026_25Sites_xinoAmbiguity_pe20_delta0_02.png}
\end{minipage}%
\begin{minipage}{.5\textwidth}  \centering
          \includegraphics[width=\linewidth]{figures/aux_input/map_z2026_25Sites_xi0.35_pe20_delta0_02.png}
\end{minipage}
\end{figure}
\end{frame}
\begin{frame}{Land allocation dynamics: productivity ambiguity}
\transduration{1}

    \begin{figure}
\centering
\begin{minipage}{.5\textwidth} \centering
  \includegraphics[width=\linewidth]{figures/aux_input/map_z2027_25Sites_xinoAmbiguity_pe20_delta0_02.png}
\end{minipage}%
\begin{minipage}{.5\textwidth}  \centering
          \includegraphics[width=\linewidth]{figures/aux_input/map_z2027_25Sites_xi0.35_pe20_delta0_02.png}
\end{minipage}
\end{figure} 
\end{frame}

\begin{frame}{Land allocation dynamics: productivity ambiguity}
\transduration{1}

    \begin{figure}
\centering
\begin{minipage}{.5\textwidth} \centering
  \includegraphics[width=\linewidth]{figures/aux_input/map_z2028_25Sites_xinoAmbiguity_pe20_delta0_02.png}
\end{minipage}%
\begin{minipage}{.5\textwidth}  \centering
          \includegraphics[width=\linewidth]{figures/aux_input/map_z2028_25Sites_xi0.35_pe20_delta0_02.png}
\end{minipage}
\end{figure}
\end{frame}

\begin{frame}{Land allocation dynamics: productivity ambiguity}
\transduration{1}

    \begin{figure}
\centering
\begin{minipage}{.5\textwidth} \centering
  \includegraphics[width=\linewidth]{figures/aux_input/map_z2029_25Sites_xinoAmbiguity_pe20_delta0_02.png}
\end{minipage}%
\begin{minipage}{.5\textwidth}  \centering
          \includegraphics[width=\linewidth]{figures/aux_input/map_z2029_25Sites_xi0.35_pe20_delta0_02.png}
\end{minipage}
\end{figure}
\end{frame}


\begin{frame}{Land allocation dynamics: productivity ambiguity}
\transduration{1}

    \begin{figure}
\centering
\begin{minipage}{.5\textwidth} \centering
  \includegraphics[width=\linewidth]{figures/aux_input/map_z2030_25Sites_xinoAmbiguity_pe20_delta0_02.png}
\end{minipage}%
\begin{minipage}{.5\textwidth}  \centering
          \includegraphics[width=\linewidth]{figures/aux_input/map_z2030_25Sites_xi0.35_pe20_delta0_02.png}
\end{minipage}
\end{figure}
\end{frame}


\begin{frame}{Land allocation dynamics: productivity ambiguity}
\transduration{1}

   \begin{figure}
\centering
\begin{minipage}{.5\textwidth} \centering
  \includegraphics[width=\linewidth]{figures/aux_input/map_z2031_25Sites_xinoAmbiguity_pe20_delta0_02.png}
\end{minipage}%
\begin{minipage}{.5\textwidth}  \centering
          \includegraphics[width=\linewidth]{figures/aux_input/map_z2031_25Sites_xi0.35_pe20_delta0_02.png}
\end{minipage}
\end{figure}
\end{frame}

\begin{frame}{Land allocation dynamics: productivity ambiguity}
\transduration{1}

   \begin{figure}
\centering
\begin{minipage}{.5\textwidth} \centering
  \includegraphics[width=\linewidth]{figures/aux_input/map_z2032_25Sites_xinoAmbiguity_pe20_delta0_02.png}
\end{minipage}%
\begin{minipage}{.5\textwidth}  \centering
          \includegraphics[width=\linewidth]{figures/aux_input/map_z2032_25Sites_xi0.35_pe20_delta0_02.png}
\end{minipage}
\end{figure}  
\end{frame}

\begin{frame}{Land allocation dynamics: productivity ambiguity}
\transduration{1}

   \begin{figure}
\centering
\begin{minipage}{.5\textwidth} \centering
  \includegraphics[width=\linewidth]{figures/aux_input/map_z2033_25Sites_xinoAmbiguity_pe20_delta0_02.png}
\end{minipage}%
\begin{minipage}{.5\textwidth}  \centering
          \includegraphics[width=\linewidth]{figures/aux_input/map_z2033_25Sites_xi0.35_pe20_delta0_02.png}
\end{minipage}
\end{figure} 
\end{frame}

\begin{frame}{Land allocation dynamics: productivity ambiguity}
\transduration{1}

    \begin{figure}
\centering
\begin{minipage}{.5\textwidth} \centering
  \includegraphics[width=\linewidth]{figures/aux_input/map_z2034_25Sites_xinoAmbiguity_pe20_delta0_02.png}
\end{minipage}%
\begin{minipage}{.5\textwidth}  \centering
          \includegraphics[width=\linewidth]{figures/aux_input/map_z2034_25Sites_xi0.35_pe20_delta0_02.png}
\end{minipage}
\end{figure}
\end{frame}

\begin{frame}{Land allocation dynamics: productivity ambiguity}
\transduration{1}

  \begin{figure}
\centering
\begin{minipage}{.5\textwidth} \centering
  \includegraphics[width=\linewidth]{figures/aux_input/map_z2035_25Sites_xinoAmbiguity_pe20_delta0_02.png}
\end{minipage}%
\begin{minipage}{.5\textwidth}  \centering
          \includegraphics[width=\linewidth]{figures/aux_input/map_z2035_25Sites_xi0.35_pe20_delta0_02.png}
\end{minipage}
\end{figure}  
\end{frame}

\begin{frame}{Land allocation dynamics: productivity ambiguity}
\transduration{1}

   \begin{figure}
\centering
\begin{minipage}{.5\textwidth} \centering
  \includegraphics[width=\linewidth]{figures/aux_input/map_z2036_25Sites_xinoAmbiguity_pe20_delta0_02.png}
\end{minipage}%
\begin{minipage}{.5\textwidth}  \centering
          \includegraphics[width=\linewidth]{figures/aux_input/map_z2036_25Sites_xi0.35_pe20_delta0_02.png}
\end{minipage}
\end{figure}
\end{frame}

\begin{frame}{Land allocation dynamics: productivity ambiguity}
\transduration{1}

   \begin{figure}
\centering
\begin{minipage}{.5\textwidth} \centering
  \includegraphics[width=\linewidth]{figures/aux_input/map_z2037_25Sites_xinoAmbiguity_pe20_delta0_02.png}
\end{minipage}%
\begin{minipage}{.5\textwidth}  \centering
          \includegraphics[width=\linewidth]{figures/aux_input/map_z2037_25Sites_xi0.35_pe20_delta0_02.png}
\end{minipage}
\end{figure}
\end{frame}
\end{comment}

\begin{comment}
\begin{frame}{Year of maximal reforestation in site $i$ ($b=15$)}

\begin{figure}
\centering
\begin{minipage}{.5\textwidth} 
    \centering
    \includegraphics[width=0.9\linewidth]{figures/aux_input/map_78site_unc_spike_b15.png}
    \caption*{ambiguity}
\end{minipage}%
\begin{minipage}{.5\textwidth}
    \centering
    \includegraphics[width=0.9\linewidth]{figures/aux_input/map_78site_det_spike_b15.png}
    \caption*{no ambiguity}
\end{minipage}
\end{figure}

\end{frame} 
\end{comment}

\begin{frame}{Year of maximal re/deforestation}
\begin{figure}[H]
\centering
\begin{minipage}{.5\textwidth} 
    \centering
    \includegraphics[width=\linewidth]{figures/decision_histogram_b_0_adjust.png}
    \vspace{-10pt}
    \caption*{(a) b=\$0}
\end{minipage}%
\begin{minipage}{.5\textwidth}
    \centering
    \includegraphics[width=\linewidth]{figures/decision_histogram_b_15_adjust.png}
    \vspace{-10pt}
    \caption*{(b) b=\$15}
\end{minipage}
\vspace{-10pt}
%\caption{Histogram of the years in which one of the two controls is maximal, $P^{ee} = 6.6$ for ambiguity neutral and $P^{ee} = 4.5$ for ambiguity averse. For \( b = \$0 \), 12.4\% of the sites remain unchanged under ambiguity neutrality, while 15.1\% do so under ambiguity aversion. For \( b = \$15 \), 2.0\% of the sites do not change under ambiguity neutrality, and 2.7\% remain unchanged under ambiguity aversion.}
\label{map_adjusted_combine}
\end{figure}
\bi
\i For \( b = \$0 \), 12\% of the sites remain unchanged under ambiguity neutrality, while 15\% do so under ambiguity aversion.
\i For \( b = \$15 \), 2\% of the sites do not change under ambiguity neutrality, and 3\% remain unchanged under ambiguity aversion.
\ei
    
\end{frame}

\begin{frame}\fmt{Planner Value Decomposition (200 years)}

\begin{table}[H]
\centering
\scalebox{0.7}{
\begin{threeparttable}
    \begin{tabular}[t]{ccccccc}
\toprule
& \multicolumn{3}{c}{ \makecell[c]{agricultural output value  ($\$$ billion) }} & \multicolumn{3}{c}{\makecell[c]{planner value ($\$$ billion) }} \\
\cmidrule[1pt](r){2-4} \cmidrule[1pt](r){5-7} 

\makecell[c]{$b$ \\ $(\$)$} &\makecell[c]{ambiguity \\ neutral } & \makecell[c]{ambiguity \\ aversion } &  \makecell[c]{percent \\ change} &\makecell[c]{ambiguity \\ neutral }&\makecell[c]{ambiguity \\ aversion }  & \makecell[c]{percent \\ change}\\
\midrule
0 & 360 & 279 & -23 & 242 & 186 & -23 \\
10 & 73 & 70 & -4 & 288 & 253 & -12 \\
15 & 37 & 37 & 0 & 361 & 326 & -10 \\
20 & 19 & 21 & 11 & 442 & 404 & -9 \\
25 & 15 & 16 & 7 & 526 & 485 & -8 \\

\bottomrule
\end{tabular}


\begin{tablenotes}[flushleft] 
\item \footnotesize \textbf{Notes}:  $P^a=\$41.11$, the average price under the stationary distribution. Forest services are calculated using baseline shadow price ($b=\$0$).
\end{tablenotes}
\end{threeparttable}}
\end{table}\end{frame}

\begin{frame}\fmt{Transfer cost (30 years)}

\begin{table}[H]
\centering
\scalebox{0.65}{
\begin{threeparttable}
    
\begin{tabular}[t]{ccccccc}
\toprule

& \multicolumn{3}{c}{ambiguity neutral} & \multicolumn{3}{c}{ambiguity aversion} \\
\cmidrule[1pt](r){2-4} \cmidrule[1pt](r){5-7} 
\makecell[c]{$b$ \\ (\$)} & \makecell[c]{net captured \\ emissions \\ (CO$_2$e Gt)} & \makecell[c]{discounted \\ net transfers \\  (\$ billion)} & \makecell[c]{  effective cost \\ (\$ per ton \\ of CO$_2$e)} & \makecell[c]{net captured \\ emissions \\ (CO$_2$e Gt)} & \makecell[c]{discounted \\ net transfers \\ (\$ billion)} & \makecell[c]{  effective cost \\ (\$ per ton \\ of CO$_2$e)} \\
\midrule



 0 & -15.9 & 0 & - & -16.8 & 0 & - \\
10 & 14.4 & 108 & 4.8  & 13.8 & 102 & 4.5 \\
 15 & 16.7 & 191 & 7.7 & 16.2 & 185 & 7.4 \\
 20 & 17.6 & 271 & 10.6 & 17.2 & 264 & 10.1 \\
 25 & 18.0 & 347 & 13.3 & 17.6 & 339 & 12.8 \\

\bottomrule
\end{tabular}

\begin{tablenotes}[flushleft] 
\item \footnotesize \textbf{Notes}: Agricultural price $P^a =  \$41.11$, which is the mean of the agricultural price in the stationary distribution.
\end{tablenotes}
\end{threeparttable}}
\end{table}\end{frame}

\begin{comment}
\begin{frame}
\fmt{Interactions across sites}
\bi
\i \cite{AAHS23}
\i  50\% of moisture in  Amazon  transmitted from trees-trees via atmospheric transportation  \i Wind direction and speed change. 
\i Contruct historical back trajectories and use this information to ``identify'' effect of Leaf Area Index shock in one pixel (25 km x 25 km) to other pixels. 
\i On average, deforestation has a ``multiplier' of 2.05.
\ei
\end{frame}
\end{comment}

\begin{comment}

\section{Interactions across sites}
\begin{frame}
\fmt{Interactions across sites}
\bi
\i \cite{AAHS23}.
\i Amazon produces large fraction of its own rainfall
\i Rainfall $\rightarrow$ trees' transpiration $\rightarrow$ recharges atmospheric humidity $\rightarrow$ humidity moves downwind $\rightarrow$ rainfall.
\i Less trees $\rightarrow$ less water. Deforestation $\rightarrow$ degradation.
\i Mapping transport of water: atmospheric trajectories 
\i Use variations in back trajectories to estimate impact of upwind Leaf Area Index (LAI) on downwind LAI
\i On average, deforestation has a ``multiplier" of 2.05.
\i Additional externality
\ei
\end{frame}
\begin{frame}{Multiplier Effect: \bf{A}: total effect of pixels; \bf{B}: total effect on pixels.}

\begin{figure}[H] 
\centering
\begin{minipage}{1\textwidth} \centering
    \includegraphics[scale=0.35]{figures/aux_input/TP effect_affect.png} 
\end{minipage}%
\end{figure}
\end{frame} 

\begin{frame}{Transboundary Cascading Effects}

\begin{figure}[H] 
\centering
\begin{minipage}{1\textwidth} \centering
    \includegraphics[scale=0.35]{figures/aux_input/TP Transboundary Cascading Effects.png} 
\end{minipage}%
\end{figure}
\bi
\i Rondônia: one of the most active frontiers of deforestation
\i 17 pixels (25km x 25km) in Rondônia, which are among the highest 50\% deforested pixels in the biome. 
\i Deforestation causes degradation as far as Bolivia
\i Deforestation multiplier of 1.87.
\ei
\end{frame} 
    
\end{comment}


\section{Conclusion}
\begin{frame}  
\fmt{Conclusions}\vspace{-12pt}
\bi
\i Posed  explicit dynamic model across heterogeneous regions in  Amazon to assess  potential adverse impact of deforestation. 
\i  Rich panel data set
\i Computational challenge because the heterogeneity of subregions requires large number of state variables and  state-constraints that bind at optimum.
\bi \i Parameter uncertainty \ei
\i With modest prices for CO2e, Brazilian Amazon would produce noticiable CO$_2$ capture.
\bi \i Compared to IPCC budget \i Compared to \cite{griscom2017natural} that identify
and quantify ``natural climate solutions'' (NCS).   \ei  \i Interactions across sites make predicted path under ``business as usual'' even more perilous. \ei \end{frame}
\begin{comment}
\begin{frame}  
\fmt{To add}
\bi \i Incentive for planner not to break contract\i Effect of Amazon deforestation on rainfall in Amazon and other regions.
\i Local tipping points
\i Other tropical forests.
\ei\end{frame}
\end{comment}



%%%% ADDITIONAL SLIDES ON HMC AND CONJUGATE PRIORS
\begin{comment}
\begin{frame}{Conjugate priors for baseline $\pi$}
\PH{OLD SLIDE}
\bi
\i To construct Bayesian posteriors for the baseline $\pi(\rho)$, we first specify a model likelihood and prior distribution.
\i We use a weighted regression model with a Gaussian error term for agricultural productivity:
\begin{equation}\label{eq:muni_likelihood}
   Y_\theta  \;=\; R_\theta\beta_\theta + \varepsilon_\theta, \quad \varepsilon_\theta \sim \mathcal{N}(0,\, \sigma_\theta^2 W_\theta ^{-1}),
\end{equation}
\i The model for $\beta_\gamma$ is analogous, with $W_\gamma = I$.

\i We assume that the priors for parameters $(\beta_\theta, \sigma_\theta^2)$ take normal-inverse gamma form:
\begin{equation*}
    \beta_\theta |  \sigma_\theta^2  \sim \mathcal{N}(m_0, \sigma_\theta^2 Q_0^{-1}), \quad \sigma_\theta^2 \sim \textrm{Inv-Gamma}(a_0,\, b_0),
\end{equation*}

\ei
\end{frame}

\end{comment}

\begin{comment}
    
\begin{frame}{Deriving posteriors for baseline $\pi$}
\PH{OLD SLIDE}
\bi

\i Then Bayesian posteriors belong to the same family of distributions, and take the form:
\begin{equation*}
    \beta_\theta |  \sigma_\theta^2  \sim \mathcal{N}(m, \sigma_\theta^2 Q^{-1}), \quad \sigma_\theta^2 \sim \textrm{Inv-Gamma}(a,\, b),
\end{equation*}
where:
\begin{align*}
    Q & = R_\theta' W_\theta R_\theta + Q_0, \cr
    m & = Q^{-1} (R_\theta' W_\theta y_\theta + Q_0m_0), \cr
    a & = a_0+\frac{n}{2}, \cr
    b & = b_0+\frac{1}{2}(y_\theta' W_\theta y_\theta + m_0' Q_0 m_0- m'Qm),
\end{align*}
\i By setting prior inputs to the limit cases $Q_0 = 0, m_0 = 0, a_0 = b_0 = 0$ we get improper uniform priors over $(\beta_\theta, \sigma_\theta^2)$ and recover familiar regression statistics for posteriors.

\ei
\end{frame}

\end{comment}
\begin{frame}[allowframebreaks]\fmt{Deriving posteriors for baseline $\pi$}
\bi \i We consider the random effects regression models for $\theta$ and $\gamma$ separately. 
\i Given observations $i=1,\ldots, N$ and groups $b=1,\ldots, B$, we  using the following specification:
\begin{equation}\label{eq:regression_general}
Y = R\beta +Z\nu +U
\end{equation}
where the observed outcome $Y$ is N-dimensional, $\beta$ is $K$ dimensional, the group indicator matrix $Z$  is $N \times B$ dimensional, and the vector of random coefficients $\nu$ is B-dimensional. 
\i 
$$
U \sim \textrm{Normal} \left(0, \frac{1}{\eta}\Sigma \right),
$$
$\Sigma$ is a  diagonal matrix. 
\i Set $W$ be such that $WW'=\Sigma^{-1}$.
\i We transform \eqref{eq:regression_general} by pre-multiplying  by $W$. 
\i $$
WY | W,R,Z,\beta,\nu,\eta \sim \textrm{Normal}\left( WR \beta+WZ\nu, \  \frac{1}{\eta} I \right).
$$
\i Impose normal priors on both vectors of coefficients including the random effects, with precision parameters $\eta$ and $\zeta$, respectively:
\[
\begin{bmatrix} \beta \cr \nu \end{bmatrix} \mid \eta, \zeta  \sim \textrm{Normal} \left(
\begin{bmatrix} 0 \cr 0 \end{bmatrix}, \begin{bmatrix} \frac 1 \eta \Omega^{-1} & 0 \cr 0 &
I \end{bmatrix} \right)
\]
\i Posterior is  
\[
\begin{bmatrix} \beta \cr \nu \end{bmatrix} \mid R, Z, \eta, \zeta  
\] 
as normal with mean
\begin{equation} \label{posterior_mean}
\begin{bmatrix} \eta R'\Sigma^{-1} R  & \eta R'\Sigma^{-1} Z
\cr \eta Z' \Sigma^{1} R & \eta Z'\Sigma^{-1} Z + \zeta I \end{bmatrix}^{-1} \begin{bmatrix} \eta R' \Sigma^{-1} Y \cr \eta Z' \Sigma^{-1}Y \end{bmatrix}   
\end{equation}
and covariance matrix 
\begin{equation} \label{posterior_covariance}
 \begin{bmatrix} \eta R'\Sigma^{-1} R  & \eta R'\Sigma^{-1} Z
\cr \eta Z' \Sigma^{1} R & \eta Z'\Sigma^{-1} Z + \zeta I \end{bmatrix}^{-1} 
\end{equation}
where we take limits as $\Omega$ tends to zero.  

Finally impose improper uniform priors on $\log(\eta)$ and $\log(\zeta)$ and thus the implied  improper prior for $\eta$ and $\zeta$ is:
\[
\eta^{-1} \zeta^{-1} 
\]
Incorporating this, gives the implied logarithm of the density for $\eta$ and $\zeta.$
\begin{equation} \label{marginal_posterior}
\begin{split}
& \frac {N-2}  2 \log\eta + \frac K 2 \log\eta + \frac {B - 2}  2 \log \zeta
- {\frac 1 2} \log  \det  \begin{bmatrix} \eta R'\Sigma^{-1} R   & \eta  R' \Sigma^{-1} Z \cr \eta Z'\Sigma^{-1} X & \eta Z'Z + \zeta  {\mathbb I} \end{bmatrix} \cr
& - \frac \eta 2 R'\Sigma^{-1} R + \frac 1 2
\begin{bmatrix} \eta Y' \Sigma^{-1} R   \eta Y' \Sigma^{-1} Z \end{bmatrix}
\begin{bmatrix} \eta R'\Sigma^{-1}R    \eta  R'\Sigma^{-1} Z \cr \eta Z'\Sigma^{-1} X  \eta Z'\Sigma^{-1}Z + \zeta  {\mathbb I} \end{bmatrix}^{-1} \begin{bmatrix} \eta R'\Sigma^{-1}Y \cr  \eta Z'\Sigma^{-1} Y \end{bmatrix}
\end{split}
\end{equation}


One way we produce the baseline posterior probabilities for the parameters of interest is first to random sample of  $(\eta,\zeta)$'s  numerically using the exponential the expression in \eqref{marginal_posterior}.   Next draw a corresponding sample of regressions coefficients using a normal distribution with mean \eqref{posterior_mean} and covariance matrix \eqref{posterior_covariance}.   Appendix \ref{sec:HMC}, which includes an adjustment for ambiguity aversion, provides a second way in the limiting case in which $\xi = \infty.$







In the following tables, we present quantiles for the posterior distributions described above:

\begin{table}[H]
\caption{Quantiles for $\theta$ posterior estimation}
\vspace{-12pt}
\centering
\scalebox{1}{
\begin{threeparttable}
    \begin{tabular}{ccccccccc}
\hline
 & $\beta_0^{\theta}$ & $\beta_1^{\theta}$ & $\beta_2^{\theta}$ & $\beta_3^{\theta}$ & $\beta_4^{\theta}$ & $\beta_5^{\theta}$ & $\beta_6^{\theta}$ & $\beta_7^{\theta}$ \\
\hline
10\% & 3.941 & -0.667 & -0.534 & 0.547 & -4.685 & -0.252 & -0.099 & 0.416 \\
50\% & 4.021 & -0.493 & -0.378 & 2.482 & -2.585 & -0.187 & -0.056 & 0.458 \\
90\% &  4.098 & -0.309 & -0.212 & 4.574 & -0.632 & -0.123 & -0.013 & 0.499 \\

\hline
\end{tabular}
\end{threeparttable}}
\label{tab:theta_estimation}
\end{table}

\begin{table}[H]
\caption{Quantiles for $\gamma$ posterior estimation}
\vspace{-12pt}
\centering
\scalebox{1}{
\begin{threeparttable}
    \begin{tabular}{ccccccc}
\hline
 & $\beta_0^{\gamma}$ & $\beta_1^{\gamma}$ & $\beta_2^{\gamma}$ & $\beta_3^{\gamma}$ & $\beta_4^{\gamma}$ & $\beta_5^{\gamma}$ \\
\hline
10\% &  6.205 & -0.011 & -0.082 & 0.089 & -0.195 & 0.124 \\
50\% & 6.233 & 0.002  & -0.073 & 0.112 & -0.170 & 0.147 \\
90\% &  6.264 & 0.015  & -0.064 & 0.136 & -0.144 & 0.169 \\
\hline
\end{tabular}
\end{threeparttable}}

\label{tab:gamma_estimation}
\end{table}

\begin{table}[H]
\caption{Quantiles for $\sigma$ posterior estimation}
\vspace{-12pt}
\centering
\scalebox{1}{
\begin{threeparttable}
    \begin{tabular}{ccccc}
\hline
 & $\sigma_{\eta}^{\gamma}$ & $\sigma_{\zeta}^{\gamma}$ & $\sigma_{\eta}^{\theta}$ & $\sigma_{\zeta}^{\theta}$ \\
\hline
10\% &  0.118 & 0.187 & 0.395 & 0.254  \\
50\% & 0.121 & 0.208  & 0.413 & 0.326  \\
90\% &  0.125 & 0.233  & 0.432 & 0.414\\
\hline
\end{tabular}
\end{threeparttable}}

\label{tab:sigma_quantiles}
\end{table}

In Figure \ref{fig:bayesian_R2} we report Bayesian R-squared for the regression equations, computed following the methodology of \cite{gelman2019r}.

\begin{figure}[H]
\centering
\begin{minipage}{.5\textwidth} 
    \centering
    \includegraphics[width=1.0\linewidth]{figures/aux_input/bayesian_r2_gamma_1043.png}
    % \caption*{(a) $\gamma$}
\end{minipage}%
\begin{minipage}{.5\textwidth}
    \centering
    \includegraphics[width=1.0\linewidth]{figures/aux_input/bayesian_r2_theta_1043.png}
    % \caption*{(b) $\theta$}
\end{minipage}
\vspace{-10pt}
\caption{The densities of Bayesian $R^2$ for the regression equations we use to recover the  $\gamma^i$'s  and $\theta^i$.} 
\label{fig:bayesian_R2}
\end{figure}

In what follows, let $\log {\mathcal L}(Y \mid R, Z, \rho)$ denote the log-likelihood function, $\pi (\rho)$ the baseline prior density.  

    \bi\i Implement ``similar''  random effects regression models for $\theta$ and $\gamma$ separately. 
    \i Given observations $i=1,\ldots, N$ and groups $b=1,\ldots, B$ set:
\begin{equation}\label{eq:regression_general}
Y = R\beta +Z\nu +U
\end{equation}
Where  observed outcome $Y$ is N-dimensional,  vector of coefficients $\beta$ is $K$ dimensional,  group indicator matrix $Z$  is $N \times B$ dimensional, and vector of random coefficients $\nu$ is B-dimensional. \i
$$
U \sim \textrm{Normal} \left(0, \frac{1}{\eta}\Sigma \right),
$$
 $\Sigma$ diagonal matrix.
 \i Let t $W$ be such that $WW'=\Sigma^{-1}$.
 \i Thus 
$$
WY | W,R,Z,\beta,\nu,\eta \sim \textrm{Normal}\left( WR \beta+WZ\nu, \  \frac{1}{\eta} I \right).
$$

$$
\beta \mid \eta, \zeta  \sim \textrm{Normal} \left(0, \frac 1 \eta \Omega^{-1}\right), \quad \nu \mid \eta, \zeta \sim \textrm{Normal} \left( 0, \frac 1 \zeta I \right).
$$

\i Impose improper uniform priors on $\log(\eta)$ and $\log(\zeta)$, and take  limit as $\Omega\to 0$. 
\i We can then sample from the posterior density, which is proportional to the model-implied likelihood and our choice of priors:
\begin{equation}
\label{baseline-distribution}
    {\overline \pi}(\rho \mid W, R, Y, Z) \ \propto \ {\mathcal L}(Y|W, R, Z, \rho) \pi(\rho ),
\end{equation}
\i For $\theta$ , observations are municipalities, groups are Amazonian water basins, $y_i$ is the log-value of cattle/hectare, and weights are the 2017 pasture area in each municipality. 
For  $\gamma$ , observations are fine (1043) sites, groups are coarse sites, $y_i$ is the average log-CO$_2$ density in each site $i$, and $W = I$. 
\ei \end{frame}


% We assume diffuse priors over the parameter vector $(\beta, \eta, \nu, \tau)$ and sample posterior distributions using the following Gibbs sampling routine:

% \begin{itemize}
%     \item Index algorithm iterations by $t$ and initialize parameters $(\beta_{(0)}, \eta_{(0)}, \nu_{(0)}, \tau_{(0)})$.
%     \item Transform the regression data and group shocks using $W$:
%         $$
%         \tilde{y} = Wy, \quad \tilde{R} = WR,\quad \tilde{\nu}_{(0)} = WG\nu_{(0)}, 
%         $$ where $G$ is a $N\times B$ matrix such that $G_{ib} = \mathbb{1}\{i\in b\}$
%     \item  For $t=1,\ldots,T$, repeat the following steps:
% \end{itemize}
% \begin{enumerate}
%     \item Sample $\beta_{(t+1)} \sim \mathcal{N}(m_{(t+1)}, Q^{-1}_{(t+1)})$, where
%     \begin{align*}
%         Q_{(t+1)} = \eta_{(t)}(\tilde{R}'\tilde{R}), \quad
%         m_{(t+1)} = Q^{-1}_{(t+1)}\eta_{(t)}\tilde{R}'(\tilde{y}-\tilde{\nu}_{(t)})
%     \end{align*}
%     \item Sample $\eta_{(t+1)} \sim \text{Gamma}(a_{(t+1)},b_{(t+1)})$ where
%     \begin{align*}
%         a_{(t+1)} = \frac{N}{2}, \quad b_{(t+1)} = \frac{\hat{\varepsilon}_{(t+1)}' \hat{\varepsilon}_{(t+1)}}{2}, \quad 
%         \hat{\varepsilon}_{(t+1)} = (\tilde{y}-\tilde{R}\beta_{(t+1)} - \tilde{\nu}_{(t)})
%     \end{align*}
%     \item Sample $\tau_{(t+1)} \sim \text{Gamma}(c_{(t+1)}, d_{(t+1)})$ where 
%     $$
%     c_{(t+1)} = \frac{B}{2}, \text{ and } d_{(t+1)} = \frac{\nu_{(t)}'\nu_{(t)}}{2}
%     $$
%     \item For each group $b = 1,\ldots, B$, sample 
%     $
%     \nu_{b,(t+1)} \sim \mathcal{N}(s_{b,(t+1)},p^{-1}_{b,(t+1)}),
%     $ where
%     \begin{align*}
%         s_{b,(t+1)} &= \frac{\eta_{(t+1)} \overline{\nu_b}}{p_{(t+1)}},\quad p_{b,(t+1)} = n_b\eta_{(t+1)}  + \tau_{(t+1)}\\[10pt]
%         \overline{\nu}_b &= {\frac 1 {n_b}} \sum_{i\in b} W_{ii}^2\left(\tilde{y_{i}} - \tilde{R_i}\cdot \beta_{(t+1)}\right), \quad n_b = \sum_{i\in b} W^2_{ii} 
%     \end{align*}
% \end{enumerate}


% The steps in this routine make up a transition for a Markov process for $(\beta_{(t)}, \eta_{(t)}, \nu_{(y)}, \tau_{(t)})$. We run the algorithm for $T=500,000$ iterations, discarding the first 300 thousand as ``burn-in" samples to allow convergence onto the stationary distribution.
\begin{comment}
\begin{frame}[allowframebreaks]\fmt{Deriving posteriors for baseline $\pi$}
    \bi
    \i Given observations $i=1,\ldots, N$ and groups $b=1,\ldots, B$, we use a weighted regression model with group-specific random effects:
    
    \begin{equation}\label{eq:muni_likelihood}
    y_{i} = R_{i}\cdot \beta + \nu_{b[i]} + \varepsilon_{i}, 
\end{equation}
    \i Group-specific shocks are $\text{i.i.d.}$ and homoscedastic, with $\text{Var}(\nu_b)=\tau^{-1}$.
    \i Common shocks are i.i.d but heteroscedastic and have $\text{Cov}(\varepsilon) = \eta^{-1}\Omega$, with $\Omega$ diagonal.
    \i Let W be a matrix of weights such that $WW' = \Omega$.
    \i For $\theta$ case: observations are municipalities, groups are Amazonian water basins, $y_i$ is the log-value of cattle/hectare, and weights are the 2017 pasture area in each municipality.
    \i For $\gamma$ case: observations are fine sites, groups are coarse sites, $y_i$ is the average log-CO$_2$ density in each site $i$, and $W = I$.
    %\i We model coefficient parameters as having a conditional normal likelihood, and precision parameters as having a conditional gamma likelihood.
    \i The model is parameterized by $(\beta, \eta, \nu, \tau)$
    \bi
    \i Coefficients $\beta$ are a vector of size $K\times1$, where $K$ is the number of geographical variables included in $R_i$.
    \i Group-specific shocks $\nu$ are a vector of size $B\times1$.
    \i Precision parameters $\eta$ and $\tau$ are scalars.
    \ei
    \i If pairs $(\beta, \eta)$ and $(\nu, \tau)$ have conditional normal-gamma priors, then conditional posteriors will also be normal-gamma by conjugacy.
    \i We assume diffuse priors over parameters and use a Gibbs sampling routine to obtain posteriors distributions.
    \In Gibbs sampling is a Markov chain Monte Carlo (MCMC) algorithm for sampling from a specified multivariate probability distribution when  sampling from the joint distribution is hard, but sampling from the conditional distribution is easier.
    \i The steps in this routine make up a transition for a Markov process for $(\beta, \eta, \nu, \tau)$.
    \bi
    \i We run the algorithm for $T=500$k iterations, discarding the first 300k as ``burn-in" samples to allow convergence onto the stationary distribution.
    \ei
\ei
\end{frame}

\begin{frame}[allowframebreaks]\fmt{Gibbs sampling}
    \begin{itemize}
        \item Transform the regression data and group shocks using $W$:
        $$
        \tilde{y} = Wy, \quad \tilde{R} = WR,\quad \tilde{\nu}_{(t)} = WG\nu_{(t)}, 
        $$ where $G$ is a $N\times B$ matrix such that $G_{ib} = \mathbb{1}\{i\in b\}$
        
        %\item Also, for notational convenience, let $\widetilde{\nu_{(t)}} = WG\nu_{(t)}$, where $G$ is a $N\times B$ matrix such that $G_{ib} = \mathbb{1}\{i\in b\}$

        \item Index algorithm iterations by $t$, initialize parameters $(\beta_{(0)}, \eta_{(0)}, \nu_{(0)}, \tau_{(0)})$, and do the following for $t=1,\ldots,T$:
    \end{itemize}
    
    \begin{enumerate}
        \item Sample $\beta_{(t+1)} \sim \mathcal{N}(m_{(t+1)}, Q^{-1}_{(t+1)})$, where
        \begin{align*}
        Q_{(t+1)} = \eta_{(t)}(\tilde{R}'\tilde{R}), \quad
        m_{(t+1)} = Q^{-1}_{(t+1)}\eta_{(t)}\tilde{R}'(\tilde{y}-\tilde{\nu}_{(t)})
        \end{align*}
            
        \item Sample $\eta_{(t+1)} \sim \text{Gamma}(a_{(t+1)},b_{(t+1)})$ where
            \begin{align*}
            a_{(t+1)} &= \frac{N}{2} \text{ , } b_{(t+1)} = \frac{\hat{\varepsilon}_{(t+1)}' \hat{\varepsilon}_{(t+1)}}{2},\\[10pt]
            \hat{\varepsilon}_{(t+1)} &= (\tilde{y}-\tilde{R}\beta_{(t+1)} - \tilde{\nu}_{(t)})
            \end{align*}

        \item Sample $\tau_{(t+1)} \sim \text{Gamma}(c_{(t+1)}, d_{(t+1)})$ where 
    $$
    c_{(t+1)} = \frac{B}{2}, \text{ and } d_{(t+1)} = \frac{\nu_{(t)}'\nu_{(t)}}{2}
    $$

    \item For each group $b = 1,\ldots, B$, sample 
    $$
    \nu_{b,(t+1)} \sim \mathcal{N}(s_{b,(t+1)},p^{-1}_{b,(t+1)}),
    $$ where
    \begin{align*}
        s_{b,(t+1)} &= \frac{\eta_{(t+1)} \overline{\nu_b}}{p_{(t+1)}},\quad p_{b,(t+1)} = n_b\eta_{(t+1)}  + \tau_{(t+1)} \\[10pt]
        \overline{\nu}_b &= {\frac 1 {n_b}} \sum_{i\in b} W_{ii}^2\left(\tilde{y_{i}} - \tilde{R_i}\cdot \beta_{(t+1)}\right), \quad n_b = \sum_{i\in b} W^2_{ii} 
    \end{align*}

    \end{enumerate}
\end{frame}

\end{comment}
\begin{frame}
\fmt{HMC - Setup}
\bi
\i Define vector $\rho := (\beta, \nu) \in \mathbb{R}^n$. We want to sample from:
\begin{equation}
\label{eq:target_distribution}
g_d(\rho)d\pi(\rho) \; \propto \; \exp\left[-\frac 1 \xi  f(d, \rho)\right]d\pi(\rho)
\end{equation}
\i By taking logs, define potential energy term:

$$
\mathcal{U}(\rho):= \frac 1 \xi  f(d, \rho) - \log d\pi(\rho) 
$$ 
\i Consider an auxiliary ``momentum'' variable  $\omega\in R^n$ where $\omega \sim \mathcal{N}(0, M)$. Then the system Hamiltonian is given by:
$$
\mathcal{H}(\rho, \omega) :=  \mathcal{U}(\rho) + \frac{1}{2} \omega'M^{-1}\omega
$$
\ei
\end{frame}

\begin{frame} \fmt{Hamiltonian Dynamics}
\bi
\i We can simulate the motion of a point in our sample space using Hamilton's equations:

\begin{equation}
\label{eq:hamiltonian_dynamics}
    \frac{d \rho}{d t} = \frac{\partial \mathcal{H}}{\partial \omega}, \quad \frac{d \omega}{d t}  = -\frac{\partial \mathcal{H}}{\partial \rho}
\end{equation}

\i Reversibility, invariance and volume preserving.
\i $\mathcal{H}$, ``energy function'',   defines  canonical distribution. %with probability density function, $$ P( \beta, \mu)=\frac 1 Z \exp(-U(\beta))\exp(-K(\mu))$$
with density
$$
P(\rho, \omega) = \frac{1}{Z} \exp\left(-\mathcal U(\rho)\right)\exp\left(-\frac{1}{2} \omega'M^{-1}\omega\right)
$$ 
\i Invariance $\Rightarrow$ dynamics moves in subset of constant prob.   \ei
\end{frame}

\begin{frame}
\fmt{HMC - The algorithm}
\vspace{-10pt}
\begin{enumerate}
    \item Initialize $\rho_{(0)}$ and sample momentum $\omega_{(0)} \sim N(0,M)$.

    \item Generate a state proposal $(\tilde{\rho}_{(0)}, \tilde \omega_{(0)})$ by evolving its position according to dynamic system in \eqref{eq:hamiltonian_dynamics}, using the leapfrog integrator with step size $\epsilon$ for $L$ steps.
\bi \i Preserves canonical probability\ei
    \item Perform a Metropolis test for the state update
          $(\rho_{(1)}, \omega_{(1)}) \leftarrow (\tilde{\rho}_{(0)}, \tilde \omega_{(0)})$, with the
          acceptance probability:
          \begin{equation*}
              \min \left\{1, \ \exp \left( \mathcal{H}(\rho_{(0)}, \omega_{(0)}) - \mathcal{H}(\tilde{\rho}_{(0)}, \tilde \omega_{(0)}) \right) \right\}
          \end{equation*}

    \item Repeat steps 2-4 until 1000 samples (plus 500 burn-in samples) are drawn.
\end{enumerate}\\[10pt]

Note: We use the Stan software implementation, which uses the NUTS algorithm, an adaptive variant of HMC.
\end{frame}

%\begin{frame}
%\fmt{HMC - The algorithm}

%\begin{enumerate}
   % \item 
%\end{enumerate}
%\end{frame}

\begin{frame}
\fmt{HMC - Solving for $g_d(\rho)$}
We then iterate between solving the planner's problem for $d$ and sampling $\rho$ as follows:

\begin{enumerate}
    \item Initialize $\varphi_{(0)}$ as the transformed mean under $\pi$.
    \item Solve the planner's problem for $d_{(0)}$ using $\varphi_{(0)}$.
    \item Sample $\{ \rho_{(s)} \}^{4000}_{s=1}$ from expression (\ref{eq:target_distribution}) by running HMC simultaneously across 4 independent Markov chains, taking 1000 samples and 500 burn-in samples per chain.
    \item Transform samples $\{ \beta_{(s)} \}^{4000}_{s=1}$ back into the $\varphi$ space, compute $\bar{\varphi}$ as the mean across samples, and update $\varphi$ using $\varphi_{(t+1)} := w \bar{\varphi} + (1-w) \varphi_{(t)}$, with $w=0.25$.
    
    \item Repeat steps 2- 4 until $||\varphi_{(t+1)} - \varphi_{(t)} ||_{\infty} < 0.001$.
\end{enumerate}
\end{frame}

\section{References}
\begin{frame}[allowframebreaks]\fmt{References}
%\bibliographystyle{apalike}
%\bibliography{academic_papers.bib}
\bibliography{Clect,Amazon}
\end{frame}

\end{document}
